\section*{Bab \ref{ch:g10:measuring-the-earth} --- Mengukur Bumi}

\begin{solution}{exo:g10:measuring-the-earth:1}
Kapal menghilang mulai dari lambungnya: di atas laut yang datar kapal hanya
akan mengecil, dan seluruhnya tetap terlihat. Bayang-bayang Bumi di Bulan saat
gerhana selalu berupa busur lingkaran: piringan datar akan menghasilkan
bayang-bayang lonjong bila terlihat dari samping. Bila kita bepergian ke
utara, Bintang Utara naik di atas cakrawala: pada Bumi yang datar tingginya
hampir tidak berubah.
\end{solution}

\begin{solution}{exo:g10:measuring-the-earth:2}
\omterm{def:g10:measuring-the-earth:geodesy}{Geodesi} mengukur bentuk dan ukuran Bumi serta posisi titik-titik di
permukaannya, beserta perubahannya dari waktu ke waktu. \omterm{def:g10:measuring-the-earth:triangulation}{Triangulasi}
menentukan letak titik-titik dengan rantai segitiga yang semua sudutnya
diukur; hanya satu panjang, yaitu garis dasar, yang diukur di tanah.
\end{solution}

\begin{solution}{exo:g10:measuring-the-earth:3}
$C = 50 \times 5000 = 250\,000$ stadion; $250\,000 \times \qty{157.2}{m} =
\qty{39300}{km}$.
\end{solution}

\begin{solution}{exo:g10:measuring-the-earth:4}
$111/60 \approx \qty{1.85}{km}$. Seperempat meridian: $90 \times 111 \approx
\qty{10000}{km}$.
\end{solution}

\begin{solution}{exo:g10:measuring-the-earth:5}
$a - b = 6378.137 - 6356.752 = \qty{21.385}{km}$;
$f = 21.385/6378.137 \approx 0.003353 \approx 1/298$.
\end{solution}

\begin{solution}{exo:g10:measuring-the-earth:6}
Selatan: $\tan\theta_1 = 0.364$, $\theta_1 \approx \qty{20.0}{\degree}$. Utara:
$\tan\theta_2 = 0.456$, $\theta_2 \approx \qty{24.5}{\degree}$.
$\alpha = \qty{4.5}{\degree}$, jadi $C = (360/4.5) \times 500 = 80 \times 500 =
\qty{40000}{km}$.
\end{solution}

\begin{solution}{exo:g10:measuring-the-earth:7}
$C = 180 - 62 - 71 = \qty{47}{\degree}$.
$AC = 10.00 \sin 71^\circ/\sin 47^\circ = 10.00 \times 0.9455/0.7314 \approx
\qty{12.93}{km}$; $BC = 10.00 \sin 62^\circ/\sin 47^\circ \approx
\qty{12.07}{km}$.
\end{solution}

\begin{solution}{exo:g10:measuring-the-earth:8}
Perubahan bujurnya $10 + 20 = \qty{30}{\degree}$. Pada
\qty{60}{\degree} LU satu derajat panjangnya $111.3 \times 0.5 = \qty{55.65}{km}$, jadi
jaraknya $30 \times 55.65 \approx \qty{1670}{km}$. Di khatulistiwa:
$30 \times 111.3 \approx \qty{3340}{km}$, dua kali lebih jauh.
\end{solution}

\begin{solution}{exo:g10:measuring-the-earth:9}
Garis pandang menyinggung laut di $T$: garis singgung lingkaran tegak lurus
pada jari-jari di titik singgungnya, jadi sudut di $T$ siku-siku. Menurut
Pythagoras: $(R+h)^2 = R^2 + d^2$, jadi $d^2 = 2Rh + h^2$; $h$ sangat kecil
dibandingkan dengan $R$, jadi $d \approx \sqrt{2Rh} = \sqrt{2 \times
6\,371\,000 \times 2} \approx \qty{5050}{m}$, sekitar \qty{5}{km}.
\end{solution}

\begin{solution}{exo:g10:measuring-the-earth:10}
Kurva paling curam di lintang tengah, sekitar \qty{30}{\degree} sampai
\qty{60}{\degree}, dan landai di dekat khatulistiwa dan kutub. Pada Bumi yang
pepat permukaan kurang melengkung di dekat kutub; makin kurang lengkung
permukaannya, makin jauh jarak yang harus ditempuh agar garis vertikal (dan
dengan demikian lintang) berputar satu derajat.
\end{solution}

\begin{solution}{exo:g10:measuring-the-earth:11}
$1072.2/9.646 \approx \qty{111.15}{km}$ per derajat; $90 \times 111.15
\approx \qty{10004}{km}$, \qty{2}{km} lebih panjang daripada seperempat
meridian yang sebenarnya. Busur itu berpusat di dekat \qty{46}{\degree}, tempat
satu derajat sedikit lebih panjang daripada rata-rata derajat dari
khatulistiwa sampai kutub ($10\,002.0/90 \approx \qty{111.13}{km}$).
\end{solution}

\begin{solution}{exo:g10:measuring-the-earth:12}
Dengan \qty{7.1}{\degree}: $C = (360/7.1) \times 5000 \approx 253\,500$ stadion,
bukan $250\,000$: sekitar 1.4\,\% terlalu besar (sudut yang lebih kecil untuk
jarak yang sama berarti Bumi yang lebih besar). Jaraknya, yang diukur dengan
langkah para pejalan kaki di sepanjang Sungai Nil, jauh kurang teliti daripada
sudutnya.
\end{solution}

\begin{solution}{exo:g10:measuring-the-earth:13}
Pada $ABC$: $C = \qty{50}{\degree}$, $AC = 6.000 \sin 75^\circ/\sin 50^\circ
\approx \qty{7.566}{km}$. Pada $ACD$: sudut di $D$ adalah $180 - 50 - 68 =
\qty{62}{\degree}$, dan $AD = 7.566 \sin 68^\circ/\sin 62^\circ \approx
\qty{7.945}{km}$.
\end{solution}

\begin{solution}{exo:g10:measuring-the-earth:14}
$\cos\varphi = 1/2$ memberikan $\varphi = \qty{60}{\degree}$; $\cos\varphi = 0.1$
memberikan $\varphi \approx \qty{84.3}{\degree}$. Satu derajat bujur pada
lintang $\varphi$ sama dengan satu derajat di khatulistiwa dikalikan
$\cos\varphi$: perbandingannya hanya bergantung pada sudut, bukan pada
jari-jari.
\end{solution}

\begin{solution}{exo:g10:measuring-the-earth:15}
Di khatulistiwa: $111.195 - 110.574 = \qty{0.621}{km} = \qty{621}{m}$; di
kutub: $111.694 - 111.195 = \qty{0.499}{km} = \qty{499}{m}$. Kesalahan beberapa
ratus meter per derajat, yang bertambah derajat demi derajat, jauh lebih besar
daripada satu meter: posisi yang dihitung dari lintang dan bujur harus memakai
elipsoid.
\end{solution}

\begin{solution}{pb:g10:measuring-the-earth:1}
\textbf{1.} $31.20 - 24.09 = \qty{7.11}{\degree}$.

\textbf{2.} $7.11/360 \approx 0.0198 = 1/50.6$, sangat dekat dengan
$1/50$ milik Eratosthenes.

\textbf{3.} $C = 788 \times 360/7.11 \approx \qty{39900}{km}$.

\textbf{4.} Busur meridian di antara kedua paralel: selisih sudut tengah hari
adalah selisih lintang, yang hanya mengukur jarak utara--selatan kedua tempat;
komponen timur--barat jaraknya tidak berperan.

\textbf{5.} Pada tengah hari setempat Matahari berada di bidang meridian, yang
sama untuk kedua tempat; pada hari yang sama Matahari menempati posisi yang
sama di antara bintang-bintang, sehingga kedua sudut berbeda hanya karena
kelengkungan Bumi.

\textbf{6.} $180 - 58 - 66 = \qty{56}{\degree}$.

\textbf{7.} $AC = 12.000 \sin 66^\circ/\sin 56^\circ = 12.000 \times
0.9135/0.8290 \approx \qty{13.223}{km}$.

\textbf{8.} $BC = 12.000 \sin 58^\circ/\sin 56^\circ \approx \qty{12.275}{km}$.

\textbf{9.} Sudut di $D$ adalah $180 - 47 - 75 = \qty{58}{\degree}$;
$AD = 13.223 \sin 75^\circ/\sin 58^\circ \approx \qty{15.061}{km}$.

\textbf{10.} $15.061 \cos 8^\circ \approx \qty{14.914}{km}$.

\textbf{11.} $14.914/111.13 \approx \qty{0.1342}{\degree}$, yaitu
$0.1342 \times 60 \approx 8.05$ menit busur.

\textbf{12.} $111.51 - 110.57 = \qty{0.94}{km}$: \qty{940}{m}.

\textbf{13.} Di dekat kutub garis vertikal berputar lebih lambat terhadap
jarak: permukaannya kurang melengkung di sana. Bumi pepat di kutub dan
menggembung di khatulistiwa.

\textbf{14.} $a - b = \qty{21.385}{km}$; $f = 21.385/6378.137 \approx 1/298$.

\textbf{15.} Garis unting-unting tegak lurus pada permukaan elipsoid. Pada
elips, garis yang tegak lurus pada kurva melalui pusat hanya di khatulistiwa
dan di kutub; pada \qty{45}{\degree} garis itu memotong sumbu di bawah pusat.

\textbf{16.} $10\,000\,000$ meter, yaitu \qty{10000}{km}.

\textbf{17.} $10\,001\,966/10\,000\,000 \approx \qty{1.0001966}{m}$.

\textbf{18.} Terlalu pendek: sepersepuluh juta yang sebenarnya adalah
$1.0001966$ kali meter tahun 1799, jadi meter tahun 1799 kurang sekitar
$1.97 \times 10^{-4}$ dari panjangnya, 0.02\,\%.

\textbf{19.} Terlalu pendek. Batang itu sepersepuluh juta seperempat meridian
seperti yang dihitung pada tahun 1790-an; karena batang itu lebih pendek
daripada sepersepuluh juta seperempat meridian yang sebenarnya, seperempat
meridian hasil hitungan itu lebih pendek daripada yang sebenarnya.

\textbf{20.} $1.0001966 - 1 = 0.0001966$ m: meter yang pertama sekitar
\textbf{\qty{0.2}{mm} terlalu pendek} --- kira-kira tebal dua lembar kertas
pada satu meter, ketelitian yang mengagumkan untuk survei yang dibuat dengan
teleskop dan sudut melintasi seluruh negeri.
\end{solution}
